%% Chapter 10 ----------------------------------------------------------------
\chapter{Projects, Reports and Export}
\label{ch:projects}

This chapter describes how the state of an analysis is saved and recovered, and how
figures, processed data, reports and batches of files are exported.

\section{Projects}
\label{sec:projects}
\index{project}\index{efaproj@.efaproj file}

A project file (\file{.efaproj}\index{efaproj@.efaproj file}) stores the complete state of the analyzer so that work can
be resumed exactly where it was left. A project contains:
\begin{itemize}
  \item the loaded or combined data, with the list of source files and the combination
    method;
  \item the processing state: background subtraction, RFI cleaning, isolated burst;
  \item the view: display range, thresholds, units, axis scale, color map and graph
    properties;
  \item annotations, light curves and their settings, and a loaded STEREO/SWAVES panel
    (restored without downloading it again);
  \item the active processing preset, the operation log and the last time-window
    synchronization;
  \item the analysis session: maximum-intensity points, the fundamental/harmonic choice,
    the Analyzer's fit, density model, fold and \(t_0\), and band-splitting points and fits.
\end{itemize}

Projects are managed from the \menu{File} menu:
\begin{description}
  \item[Save Project] (\kbd{Ctrl+S}, or the toolbar tool) saves to the current project
    file, asking for a name the first time.
  \item[Save Project As...] (\kbd{Ctrl+Shift+S}) saves to a new file.
  \item[Open Project...] (\kbd{Ctrl+Shift+O}) opens a project. A project file can also be
    dropped onto the main window.
  \item[Recent Projects]\index{recent projects} lists the last ten projects opened or saved;
    \gui{Clear Recent Projects} empties the list, and an entry whose file is missing is
    removed with a message.
\end{description}
Before a project is opened, or new data are loaded, you are asked whether to save unsaved
changes.

\subsection{Restored analysis sessions}
\label{sec:restored-analysis}

When a project with an analysis session is opened, the \gui{Analysis Summary} panel reports
that the session has been restored. \menu{Processing > Analysis Session > Open Restored
Analysis} reopens the windows that were open when the project was saved --- the Maximum
Intensities window, the Burst Analyzer and the Type~II Band-splitting window --- with their
points, fits and settings.

\section{Autosave and recovery}
\label{sec:autosave}
\index{autosave}\index{recovery}

While data are loaded and there are unsaved changes, the application writes a recovery
snapshot every three minutes. The ten most recent snapshots are kept. If the application
did not close normally, it offers to recover the latest snapshot at the next start.
\menu{File > Recover Last Session...}\index{recovery!Recover Last Session} restores the most recent snapshot at any time.

\section{Exporting figures}
\label{sec:export-figure}
\index{export!figure}\index{OriginPro style}

\menu{File > Export As > Export Figure}\index{export!figure} (\kbd{Ctrl+E}, or the \gui{Export} toolbar tool)
saves the main plot as a PNG, PDF, EPS, SVG, TIFF or JPG file. Whatever renderer and theme
are on screen, the exported figure is drawn as a publication-style (``OriginPro-style'')
graph in light mode:
\begin{itemize}
  \item Arial type, a closed black frame with inward major and minor ticks on all four
    sides, a framed color scale, and time ticks on whole UT minutes;
  \item the current zoom, color map, display range, linear or logarithmic frequency axis,
    and the titles and font sizes from \gui{Graph Properties};
  \item every visible overlay: light curves, annotations, the GOES X-ray overlay, the ruler
    measurement and the STEREO/SWAVES panel.
\end{itemize}
The plot on screen keeps its own look. On Windows, if the proposed folder is not writable
(for example inside \file{C:\\Program Files}), you are asked to choose another folder.

The Maximum Intensities window, the Burst Analyzer, the band-splitting window, the
Multi-Station Comparison and the Solar Image Analysis tracking panel export their graphs in
the same style.

\section{Exporting FITS files}
\label{sec:export-fits}
\index{export!FITS}

\menu{File > Export As > Export to FIT}\index{export!FITS} (\kbd{Ctrl+F}, or the \gui{Export as FITS} toolbar
tool) writes the spectrum on show as a new FITS file, for downstream analysis or machine
learning. The processed data are exported if processing has been applied, and the raw data
otherwise; combined datasets are exported as one file.

Before saving, choose the data type (\texttt{BITPIX}\index{BITPIX keyword}): \gui{Auto} (the recommended type is
shown), \gui{8 (unsigned byte)}, \gui{16 (signed int16)} or \gui{32 (signed int32)}.
Integer types are offered because some FITS viewers cannot read 64-bit floating-point data.
The exported file keeps the original header, updated with:
\begin{itemize}
  \item \texttt{HISTORY} records naming the application, the plot type, the intensity
    units and the source file(s);
  \item \texttt{COMBINED}, \texttt{COMBMETH} (combination method) and \texttt{NFILES} for
    combined data;
  \item \texttt{BUNIT}, \texttt{DATAMIN} and \texttt{DATAMAX};
\end{itemize}
and an \texttt{AXIS} binary-table extension with the \texttt{FREQUENCY} and \texttt{TIME}
columns, as in standard CALLISTO files (see Appendix~\ref{app:file-formats}).

\section{Provenance reports and analysis logs}
\label{sec:provenance}
\index{provenance report}\index{analysis log}

For reproducibility, two structured reports can be exported from \menu{File > Export As}:
\begin{description}
  \item[Export Provenance Report...] writes a Markdown (\file{.md}) and a JSON
    (\file{.json}) summary with the sections \emph{App}, \emph{Data Source},
    \emph{Processing}, \emph{RFI}, \emph{Annotations}, \emph{Time Sync} and
    \emph{Operation Log}.
  \item[Export Analysis Log...] writes \file{<name>_analysis_log.csv} and
    \file{<name>_analysis_log.txt} with the Analyzer's fit parameters and the derived shock
    quantities.
\end{description}
These reports are suitable for laboratory notebooks, collaboration hand-overs and the
preparation of papers.

\section{Project report PDF}
\label{sec:project-report}
\index{project report}

\menu{File > Generate Project Report...}\index{project report} creates a consolidated PDF of the current project
or loaded dataset. A progress dialog is shown while the report is built. The report
contains the sections listed below; sections without content are omitted.
\begin{itemize}
  \item \emph{Data Source} and \emph{Selected FITS Header};
  \item \emph{Processing}, and \emph{RFI, Annotations, and Light Curves};
  \item \emph{Analysis Summary} and \emph{Type II Results}, including the fit equation (with
    \(t_0\) when it is not the file start), the density model, and band-splitting output;
  \item \emph{Figures}: the raw and background-subtracted spectra, light curves with the
    spectrum, the maximum-intensity fit, the Type~II band-splitting plots, the
    STEREO/SWAVES split view\index{STEREO/SWAVES!split view}, and any available GOES X-ray, GOES proton, Dst and Kp context
    plots;
  \item \emph{Full FITS Header Appendix}.
\end{itemize}
Every graph in the report is drawn in the same style as \gui{Export Figure}, on a white page
even in the dark theme.

\section{Batch processing}
\label{sec:batch}
\index{batch processing}

\menu{Processing > Batch Processing > Open Batch Processor} opens the \gui{Batch FIT
Processing} dialog (Figure~\ref{fig:batch}), which renders every FIT file in a folder as a
PNG image with consistent settings.

\screenshot[0.62\linewidth]{dialogs/batch_processing_dialog}{The \gui{Batch FIT Processing} dialog.}{fig:batch}{The Batch FIT Processing dialog with input and output folders selected, the Output Type, Processing Options and Aligned View groups, and the status line.}

\begin{description}
  \item[Input Folder, Output Folder] the folder of FIT files to process and the folder for
    the PNG images (one per input file, at 300~dpi; existing names are not overwritten).
  \item[Output Type] \gui{Raw} or \gui{Background Subtracted}.
  \item[Processing Options] the \gui{Background Method} (\gui{Mean}, \gui{Median} or
    \gui{Plotutil Median (dB)}\index{Plotutil}\index{background subtraction!Median (dB)}) and the \gui{Colormap}. Plotutil Median (dB) converts digits
    to dB with the CALLISTO scale before subtracting the median, and defaults to a display
    range of \(-1\) to 8~dB.
  \item[Aligned View] \gui{Use current display range} and \gui{Use current visual settings}
    apply the main window's display range and appearance; \gui{Load View Config...} applies
    a saved \file{.efaview.json}\index{efaview.json@.efaview.json file} file instead. These options give many station spectra
    identical time and frequency axes.
  \item[Start, Cancel, Close] start the run, cancel it after the current file, or close the
    dialog. A summary is shown when the run finishes, including any files that failed.
\end{description}
